% Proposed B.85: reviewed research, not included before checkpoint 15.
\statement{B.85}{Proved}{An actual source-compatible tangential Galois action with literal product-one cyclotomy}
\label{stmt:B85}

Let
\[
 Y=\{(v,w)\in(\mathbb C^*)^2:\ v,w\notin\mu _5,
       \ v\ne w,\ vw\notin\mu _5\},
 \qquad L=\{v+w=2\}\cap Y,
\]
with the actual base point $(399/200,1/200)$ and the common source
marking of B.83. Write $Y_{\mathbb Q}$ for the smooth $\mathbb Q$-scheme
defined by these equations, so $Y=Y_{\mathbb Q}(\mathbb C)$, and fix
$\overline{\mathbb Q}\subset\mathbb C$ for the etale--topological comparison. Write $G=\widehat{\pi _1(Y)}$, and let $H\leq G$
be the actual closed image of the source fundamental group. Let $J_c,J_m$
be the actual closed images of the source kernels for filling respectively
$v\in\mu _5\setminus\{1\}$ and $w\in\mu _5\setminus\{1\}$, and
put $I=J_cJ_m$. Let $b$ be the actual transported product-one meridian,
with the inverse-conjugator orientation of B.83. There exists a selected
continuous action
\[
 \rho:G_{\mathbb Q}\longrightarrow\operatorname{Aut}(G)
\]
preserving $H,J_c,J_m,I$ and satisfying, literally,
\begin{equation}\label{eq:B85-cyclotomy}
 \rho_\sigma(b)=b^{\chi(\sigma)}.
\end{equation}
The group $\operatorname{Aut}(G)$ has its natural inverse-limit topology
on the finite characteristic quotients of $G$.
Here $\chi:G_{\mathbb Q}\to\widehat{\mathbb Z}^{\times}$ is the full
cyclotomic character. The action is obtained from a rational tangential
base point and the actual common source path. Equality with the action
from the original rational point's own arithmetic section is not asserted.

For the actual source inertia $c_j$ at $w=2-\zeta _5^j$, there are
$h_{\sigma,j}\in H$ such that
\begin{equation}\label{eq:B85-source-inertia}
 \rho_\sigma(c_j)=h_{\sigma,j}
 c_{\overline\chi(\sigma)j}^{\chi(\sigma)}h_{\sigma,j}^{-1}.
\end{equation}
Indices are reduced modulo $5$. In particular, arithmetic transport
carries the inertia classes indexed by $1,4$ to those indexed by $2,3$.
The $H$-conjugators in this assertion are retained.

There is also a controlled arithmetic enlargement of B.83. Define
\[
 H_* =\overline{\langle c_1,c_4,g_1,g_2,g_5\rangle},\qquad
 T=\overline{\langle H_*,b\rangle},\qquad
 J_*=\overline{\langle\!\langle c_1,c_4\rangle\!\rangle}_{H_*},
\]
using that statement's actual based generators, and set
\[
\begin{split}
 H_{\rm ar}&=\overline{\langle \rho_\sigma(H_*):
                              \sigma\in G_{\mathbb Q}\rangle},\\
 K_{\rm ar}&=\overline{\langle \rho_\sigma(J_*):
                              \sigma\in G_{\mathbb Q}\rangle},\\
 J_{\rm ar}&=\overline{\langle\!\langle K_{\rm ar}
                              \rangle\!\rangle}_{H_{\rm ar}}.
\end{split}
\]
Then $H_{\rm ar}\leq H$, $J_{\rm ar}\leq J_c$, and all three displayed
subgroups are preserved by $\rho$ and normalized by $b$. In fact
$H_{\rm ar}$ is a proper subgroup of $H$: the six actual winding
characters of
\[
 v,\quad w,\quad vw-\zeta_5,\quad vw-\zeta_5^2,\quad
 vw-\zeta_5^3,\quad vw-\zeta_5^4
\]
give a surjection $H\to\widehat{\mathbb Z}^{6}$ that kills
$H_{\rm ar}$. They also kill $J_c$, so this properness does not decide
whether $J_{\rm ar}=J_c$.

\begin{proof}
\emph{B.85.1. An exact rational partial compactification.}
On $Y_{\mathbb Q}$ put
\[
 t=w-1,\qquad q=\frac{vw-1}{(w-1)^2},\qquad
 w=1+t,\qquad v=\frac{1+qt^2}{1+t}.
\]
Both $q$ and $t$ are units. Define polynomials
\[
\begin{aligned}
 W(t)&=\frac{(1+t)^5-1}{t},&
 V(t,q)&=\frac{(1+qt^2)^5-(1+t)^5}{t},\\
 D(t,q)&=-2+(q-1)t,&
 P(t,q)&=\frac{(1+qt^2)^5-1}{qt^2},\\
 F(t,q)&=(1+t)(1+qt^2)W(t)V(t,q)D(t,q)P(t,q).
\end{aligned}
\]
Each quotient displayed here is polynomial. The deleted-boundary equations
are equivalent to invertibility of $q,t,F$, giving the actual isomorphism
\[
 Y_{\mathbb Q}\simeq\operatorname{Spec}
       \mathbb Q[q,q^{-1},t,t^{-1},F(t,q)^{-1}].
\]
Indeed $w^5-1=tW$, $(1+t)^5(v^5-1)=tV$,
$(1+t)(v-w)=tD$, and $(vw)^5-1=qt^2P$; the factors
$1+t$ and $1+qt^2$ impose $w,v\ne0$.
Since $F(0,q)=250$, the scheme
\[
 U=\operatorname{Spec}\mathbb Q[q,q^{-1},t,F(t,q)^{-1}]
      \longrightarrow B=\mathbb G_{m,q}
\]
is a smooth relative curve with the entire section $S=\{t=0\}\simeq B$,
and $Y_{\mathbb Q}=U\setminus S$. Its relative normal line is trivialized by $t$,
so
\[
 N^\times_{S/U}\simeq\mathbb G_{m,q}\times\mathbb G_{m,n}.
\]
Geometrically, writing $x=vw-1$, the substitutions $x=tq_1$, $q_1=tq$
are the two blowups at the rational successive centers. In the finite-$q$
chart the product-one proper transform is $q=0$; the earlier exceptional
divisor is at $q=\infty$. The strict equations $v=1$ and $v=w$ become
$qt=1$ and $qt=2+t$, respectively. Thus they do not remove any point of
the finite-$q$ section $S$. On the actual source,
\[
 v=1-t,\qquad w=1+t,\qquad vw=1-t^2,\qquad q=-1.
\]
Its negative-real normal direction is the rational point
$(q,n)=(-1,-1)$ of $N^\times_{S/U}$.

\emph{B.85.2. The relative tangential functor and its arithmetic action.}
Apply the relative normal-specialization construction of
\citet[\S15.19]{BDeligne1989} to $U/B$ and the section $S/B$.
Its hypotheses hold literally: $B$ is normal, $U/B$ is a smooth relative
curve, and $S/B$ is an etale section. In characteristic zero, finite etale
covers of $Y_{\mathbb Q}$ are tame along $S$. Normalization across $S$ and leading
normal terms therefore give a functor to finite etale covers of
$N^\times_{S/U}$. Descent over $\mathbb Q$, compatibility with the
classical normal construction, and its formal interpretation are supplied
by \citet[\S\S15.17--15.19, 15.22--15.26]{BDeligne1989}.
These precise relative-curve hypotheses avoid an unstated
higher-dimensional specialization theorem.

Taking the fiber at $(-1,-1)$ yields a rational tangential fiber functor.
It determines a selected action of $G_{\mathbb Q}$ on the geometric
profinite fundamental group. On the normal torus its action on
\[
 \pi _1(N^\times_{S/U,\overline{\mathbb Q}},(-1,-1))
       \simeq\widehat{\mathbb Z}(1)^2
\]
is cyclotomic in both factors, as also follows directly from the two
families of Kummer covers. In this example the normal-torus map is even
split injective: the global unit map $(q,t):Y_{\mathbb Q}\to\mathbb G_m^2$ has
normal specialization $(q,n)\mapsto(q,n)$ and thus supplies a left inverse
on geometric fundamental groups. This only identifies the actual local
two-generator subgroup; it asserts no normality of $H$.

\emph{B.85.3. Literal comparison with the actual based product-one loop.}
The common source whisker ends at the actual corner
$w_c=1+10^{-4}+i10^{-4}$ of the stored $w=1$ lasso. Holding $t=w_c-1$
constant, take the positive conventional $q$-loop
\[
 -1\longrightarrow-i\longrightarrow1\longrightarrow i\longrightarrow-1.
\]
The certificate checks four segments by exact polynomials and Sturm sequences.
It is saved in
\path{work_DE/next_marginal_source/} as
\path{check_w1_local_q_meridian.py}, with its genuine output and JSON
record. These checks give the local braid $\sigma_{10}^2$.
Under the actual inverse-conjugator convention its fiber word is
$x_{10}^{-1}$, and the already certified common source whisker transports
it to exactly $b^{-1}$. No additional conjugator is introduced.

Transport the entire pointed source/filling diagram from the tangential
fiber functor to the old base using this same geometric source path,
extended locally along $q=-1$ to the direction $-1$. The local extension
ambiguity lies in the punctured normal $t$-disc and hence commutes with
the $q$-loop. It does not change the literal $b^{-1}$ identification.
No Galois-invariant global path is required: the chosen action is transported
by an isomorphism of geometric fiber functors. Cyclotomic action on the
positive $q$-loop gives
\[
 \rho_\sigma(b^{-1})=b^{-\chi(\sigma)},
\]
which is exactly \eqref{eq:B85-cyclotomy}. The normal source loop is
cyclotomic as well.

\emph{B.85.4. The actual source and filling kernels.}
The source and both fillings are defined over $\mathbb Q$, since the set
of four nonidentity fifth roots is Galois stable. Their tangential fiber
functors are the compatible images of the same source tangential point.
Using the same source path for every map makes the whole actual based
diagram equivariant. Consequently $\rho$ preserves the actual compact
images $H,J_c,J_m,I$; no assertion about an induced topology on a
separately defined abstract subgroup is used. Pointed source inertia
functoriality gives \eqref{eq:B85-source-inertia}, with its conjugators
in the source image $H$. Choosing $\sigma(\zeta _5)=\zeta _5^2$ gives
the stated permutation of pairs of inertia classes.

\emph{B.85.5. Closed arithmetic enlargement and its exact scope.}
B.83 proves $J_*\subset J_c$, that $b$ normalizes $J_*$, and the actual
product decomposition
\[
 T/J_*\simeq\widehat{\mathbb Z}\times\widehat F_3,
\]
in which $b$ generates the first factor and the image of $H_*$ is the
second. Its actual character $\chi_b$ is minus the conventional winding
of $q$, kills $H$, and sends $b$ to $1$. Since $J_*\subset H_*$,
\[
 H_*=\ker(\chi_b|_T).
\]
In particular, $b$ normalizes $H_*$ itself.

For every closed $b$-normalized subgroup $K$, equation
\eqref{eq:B85-cyclotomy} shows that $\rho_\sigma(K)$ is normalized by
$b^{\chi(\sigma)}$, hence by $b$: the normalizer is closed and
$\chi(\sigma)$ is a full profinite unit. Apply this to $H_*$ and $J_*$.
It follows that $H_{\rm ar}$ and $K_{\rm ar}$ are $b$-normalized.
Their definitions make them $\rho$-stable, and $K_{\rm ar}\subset
J_c\cap H_{\rm ar}$. Because $J_c$ is normal in $H$, its closed normal
closure $J_{\rm ar}$ inside $H_{\rm ar}$ remains in $J_c$. Conjugation
by $b$ preserves both $H_{\rm ar}$ and $K_{\rm ar}$, so it preserves
that closed normal closure. The same reasoning gives its $\rho$-stability.

This argument can equally be applied to the stronger controlled source
subgroup of B.84 when its hypotheses are invoked. It does not identify
the arithmetic kernel enlargement with the full source kernel. In
particular, the $H$-conjugators in \eqref{eq:B85-source-inertia} and the
other source directions cannot be discarded. Thus whole-$J_c$
normality and the original marginal kernel equality remain open.

\emph{B.85.6. An actual six-character limit on the source enlargement.}
Let $\Xi:G\to\widehat{\mathbb Z}^{6}$ be the winding homomorphism
of the six units in the statement. The last four units are defined over
$\mathbb Q(\zeta_5)$, and their set is Galois stable. Consequently
arithmetic transport permutes their characters and scales them
cyclotomically; the first two scale cyclotomically. This statement on
abelian targets is unaffected by the transported section's inner
ambiguity. In particular, $\ker\Xi$ is $\rho$-stable.

On the source $v=2-w$, these six units become
\[
 2-w,\quad w,\quad w(2-w)-\zeta_5^j\quad(1\leq j\leq4).
\]
All are nonzero at the five selected punctures
$2-\zeta_5,2-\zeta_5^4,1,\zeta_5,\zeta_5^4$.
For example $|2-\zeta_5^j|>1$ for $j\ne0$, so the product
$\zeta_5^j(2-\zeta_5^j)$ is not a fifth root of unity.
Thus $\Xi(H_*)=0$, and Galois stability gives $\Xi(H_{\rm ar})=0$.
The positive source meridians at $w=2$, at $w=0$, and at one root of
each of the four quadratics $w(2-w)-\zeta_5^j$ give the six standard
basis vectors. Indeed the quadratics have distinct simple roots, are
pairwise relatively prime, and do not vanish at $0$ or $2$.
All six meridians are actual source elements; their whiskers do not
affect winding. Hence $\Xi(H)=\widehat{\mathbb Z}^{6}$, proving
$H_{\rm ar}\ne H$.

The calculation remains valid for the stronger source subgroup of B.84:
its additional positive-root loops have zero values, its deleted Nielsen
generator has the same values as the inverse original meridian, and the
two correcting source words have zero abelianization. Finally all six
units extend over the constant-family filling $Y_c$, so $\Xi(J_c)=0$.
The detector therefore supplies no inequality between $J_{\rm ar}$
and $J_c$.

\emph{Verification record.}
The rational chart identities and $F(0,q)=250$ were checked symbolically
in \path{work_A/code/}, file
\path{check_tangential_product_one_chart.py}, and its executed output. The independent exact local-meridian certificate above
fixes the orientation and actual whisker. The six-character divisor
calculation, including the literal source-correction abelianization and
the identity local winding matrix, was executed exactly over
$\mathbb Q(\zeta_5)$ in \path{work_A/code/}, file
\path{check_tangential_arithmetic_scope.py}, with its saved output.
The cited primary passages were
opened and their hypotheses read directly; the argument constructs a
selected action on the full profinite completion, not an unverified
identification with a rational-point section.
\end{proof}
